\chapter{系统测试与结果分析}
\section{测试环境与方法}
示例测试使用 Windows 11、XeLaTeX 论文构建环境和关系数据库测试实例。测试采用黑盒功能测试与边界条件测试相结合的方式，重点验证状态转换、库存校验和权限边界。
\section{功能测试}
\begingroup\zihao{5}\setstretch{1.15}\setlength{\tabcolsep}{3pt}\renewcommand{\arraystretch}{1.25}
\begin{longtable}{|>{\RaggedRight\arraybackslash}p{0.13\textwidth}|>{\RaggedRight\arraybackslash}p{0.35\textwidth}|>{\RaggedRight\arraybackslash}p{0.32\textwidth}|}
\caption{主要功能测试用例}\label{tab:5-1}\\\hline
\textbf{编号} & \textbf{测试操作} & \textbf{预期结果} \\\hline
\endfirsthead
T01 & 学生提交库存充足的借用申请 & 申请进入待审批状态 \\\hline
T02 & 审批库存不足的申请 & 操作失败，库存不发生变化 \\\hline
T03 & 重复登记同一申请归还 & 第二次操作被拒绝 \\\hline
T04 & 学生访问管理员审批接口 & 返回无权限错误 \\\hline
\end{longtable}\endgroup
\section{测试结果分析}
测试结果表明，正常流程和边界条件均按照设计规则执行。库存不足、重复归还和越权访问能够被服务端拦截，关键失败操作不会留下错误的状态更新。后续仍需在更大规模并发和真实部署环境下补充压力测试。
